% Drafted against audit/outline.md §2. Disclosure figures and the identity
% check come from data/processed/gs_var_average_daily.csv; Goldman's own
% wording is quoted from references/sec_filings/, both hash-recorded in
% audit/evidence_freeze.csv at commit 1e6bdd74.
%
% Two corrections to the outline's prior-work summary, made at drafting time
% against the source register: Pérignon, Deng and Wang (2008) study six
% Canadian banks, not US ones, and their abstract attributes the overstatement
% they find partly to underestimated diversification, which makes them closer
% to this paper's subject than the outline implied.

\section{The disclosed diversification effect, and prior work}\label{sec:disclosed}

\subsection{What Goldman discloses, and how the effect is defined}

The definition is Goldman's own, and it is worth quoting because the whole
exercise depends on it. In the words of the FY2025 filing, the "[d]iversification
effect in the tables below represents the difference between total VaR and the
sum of the VaR for the four risk categories," and "[t]his effect arises because
the four market risk categories are not perfectly correlated"
\citep{gs10k}. The four categories are interest rates, equity prices, currency
rates and commodity prices. The measure is a one-day 95\% VaR, and the filings
state that a single model produces both the category figures and the total,
which matters here: the diversification effect is internal to one model run
rather than the residual of aggregating four separate ones.

Table~\ref{tab:disclosed} gives the series. Across the ten fiscal years the
effect has sat between 39.0\% and 45.5\% of summed category VaR, averaging
42.0\% with a standard deviation of 2.08 percentage points and no discernible
trend. The dollar figures underneath it move a great deal: summed category VaR
runs from \$89 million in 2017 to \$209 million in 2022, and total VaR from
\$54 million to \$114 million over the same two years. The ratio of the two
barely moves.

% Generated by build/make_tables.py from gs_var_average_daily.csv. Do not edit.

\begin{table}[tbp]
\centering
\caption{Goldman Sachs average daily Value-at-Risk by risk category, and the disclosed diversification effect, FY2016--FY2025}
\label{tab:disclosed}
\begin{tabular}{@{}lrrrrrrrr@{}}
\toprule
 & \multicolumn{4}{c}{Risk category} & & & & \\
\cmidrule(lr){2-5}
Year & Rates & Equities & FX & Commod. & Summed & Diversif. & Total & \% of sum \\
\midrule
2016 & 45 & 25 & 21 & 17 & 108 & -45 & 63 & 41.7 \\
2017 & 40 & 24 & 12 & 13 & 89 & -35 & 54 & 39.3 \\
2018 & 46 & 31 & 14 & 11 & 102 & -42 & 60 & 41.2 \\
2019 & 46 & 27 & 11 & 12 & 96 & -40 & 56 & 41.7 \\
2020 & 71 & 55 & 23 & 20 & 169 & -75 & 94 & 44.4 \\
2021 & 60 & 43 & 13 & 25 & 141 & -55 & 86 & 39.0 \\
2022 & 97 & 33 & 32 & 47 & 209 & -95 & 114 & 45.5 \\
2023 & 96 & 29 & 24 & 19 & 168 & -69 & 99 & 41.1 \\
2024 & 81 & 37 & 26 & 19 & 163 & -71 & 92 & 43.6 \\
2025 & 70 & 47 & 24 & 17 & 158 & -68 & 90 & 43.0 \\
\bottomrule
\end{tabular}
\notetext{In millions of US dollars, except the final column. Average daily 95\% one-day Value-at-Risk as reported in the Market Risk Management section of each fiscal year's Form 10-K, taken from the filing that first reported the year rather than from a later filing's restatement. \emph{Summed} is the sum of the four category figures and is not itself disclosed. The diversification effect is disclosed as a negative number and reproduced with its sign; the final column expresses its magnitude as a percentage of the summed categories. Total VaR equals the summed categories plus the diversification effect in every year, exactly. The percentage averages 42.0 with a standard deviation of 2.08 across the ten years.}
\end{table}


The filings are also more specific about method than the tables alone suggest,
and three of those specifics are load-bearing later. Goldman states that it
uses "historical simulations with full valuation of market factors at the
position level by simultaneously shocking the relevant market factors for that
position"; that it samples "from five years of historical data to generate the
scenarios"; and that "[t]he historical data is weighted so that the relative
importance of the data reduces over time," which "gives greater importance to
more recent observations and reflects current asset volatilities"
\citep{gs10k}. The FY2016 and FY2017 filings put the factor count at
approximately 70,000; later filings drop the number. Section~\ref{sec:method}
adopts the method, the window and the weighting scheme; the decay rate is the
one piece of the description the filings leave unquantified.

One consequence of the weighting is stated in the filings explicitly: "even if
our positions included in VaR were unchanged, our VaR would increase with
increasing market volatility and vice versa" \citep{gs10k}. The disclosed
measure is therefore expected to respond to volatility without any change in
exposure, which is the behaviour a static replication can in principle
reproduce. What a static replication cannot reproduce is a change in exposure,
and Section~\ref{sec:results} shows that the filings attribute much of the
year-over-year movement to exactly that.

\subsection{Where the number comes from}

The figures are read out of the Market Risk Management section of each Form
10-K, filed under CIK 0000886982, for fiscal years 2016 through 2025. None of
it is available as structured data: the Item 7A tables carry no XBRL tagging,
so the numbers exist only as HTML table cells, and the cell layout is not
stable across filings. Six tables are extracted from each filing where
present: average daily VaR by category, period-end VaR, the year's high and low
by category, the daily net revenue distribution, and the exception counts at
the 95\% risk-management measure and the 99\% regulatory one. Appendix A gives
the full extracted series and the accession number of every filing.

Two features of the disclosures do the work of validating that extraction.

The first is an identity. Total VaR must equal the sum of the four categories
plus the diversification effect, because that is how the effect is defined. It
is therefore a property of a correct parse rather than a finding, and it is a
sharp one: a misread digit, a dropped minus sign or a column misalignment
breaks it. Across every average-daily row extracted --- 22 rows, since each
filing restates the prior year, covering 12 distinct disclosure years from 2014
to 2025 --- the residual is exactly zero. Net revenue bucket counts likewise sum
to the stated annual totals in all fourteen year-filing pairs that report them.

The second is redundancy. Because each filing restates the preceding year,
most years appear twice under different source filings, and the two readings
agree everywhere but once. The FY2023 filing reports 2022 average daily VaR as
\$113 million total against \$210 million summed, where the FY2022 filing had
reported \$114 million against \$209 million. The filing explains the change:
during the first quarter of 2023 Goldman added the currency exposure on certain
debt and equity positions to VaR, removed certain other positions, and
conformed the prior period. Both readings satisfy the identity, which is what
identifies this as a change in the measure's scope rather than a correction of
an error. Figures used throughout this paper are as first reported, so that a
methodology change does not sit inside the time series.

Disclosure practice is not perfectly stable either. The 99\% regulatory
exception count appears in eight of the ten filings but not in the FY2022 or
FY2023 filings, which report the 95\% count alone. Nothing in this paper
depends on the missing rows, but they are a reminder that what can be
reconstructed from disclosure is determined by the filer rather than the
researcher.

\subsection{Prior work}

The premise that these disclosures carry information is
\citet{jorion2002informative}'s: reported VaR predicts the variability of
subsequent trading revenue. What followed has largely asked how good the
reported numbers are.

Two findings recur. The first is conservatism.
\citet{perignon2008overstate} study six Canadian banks on daily data from 1999
to 2005 and find systematic overstatement, two exceptions in 7,354 trading
days against 74 expected at 99\%, which they attribute partly to banks
underestimating the diversification effect when aggregating across risk
categories. \citet{perignon2010level} construct a disclosure-quality index over
1996 to 2005 and report that historical simulation is by far the most common
method in practice, and that VaR computed that way carries very little
information about future volatility. The second is that bank VaR adapts slowly.
\citet{obrien2017evaluation} find that exceedances at large US banks were
excessive and clustered during the financial crisis, that a GARCH-based
benchmark achieved lower exceedance rates without clustering, and that variance
decompositions show limited ability of bank VaR methodologies to adjust to
crisis conditions. \citet{campbell2022australia} report the same shape in
Australia, with VaR too high in quiet periods and too low in stressed ones.

\citet{perignon2010diversification} is the nearest neighbour to this paper. It
tests whether banks underaccount for diversification across broad risk
categories, on quarterly observations for four US banks running to the first
quarter of 2007, and finds no systematic underestimation: disclosed
diversification effects are close to, and often larger than, the authors' own
estimates. The level result in Section~\ref{sec:results} corroborates that on a
different bank, a different period and a different method, and it should be read
as corroboration rather than discovery. Its sample ends before the financial
crisis and is quarterly, which leaves the behaviour of a disclosed
diversification effect across regime breaks as something a longer single-dealer
series can speak to.

\citet{danielsson2016modelrisk} supply the mechanism this paper's central
result turns on. They show that model risk, the potential for different
models to give inconsistent outcomes, rises with market uncertainty:
candidate models agree in calm periods and diverge significantly in distress.
Their object is risk forecasts in general rather than a disclosed
diversification ratio, but the diversification effect is a function of a
correlation structure estimated from data, and the finding here is a reading of
their result through that one statistic. Where their models diverge from each
other, a reconstruction diverges from the thing it reconstructs.

Finally, the methodological literature this paper borrows from directly is
small and is introduced where it is used:
\citet{kupiec1995} and \citet{christoffersen1998} for the coverage tests,
\citet{bcbs1996backtesting} for the supervisory zones,
\citet{boudoukh1998hybrid} for the age-weighted quantile that implements
Goldman's disclosed weighting, and \citet{hullwhite1998} with
\citet{riskmetrics1996} for the volatility-weighted variant.
